\chapter{Arithmetic Funtions}

\begin{definition}
    An \textbf{arithmetic function}, is a function $f: \mathbb{N} \rightarrow \mathbb{C}$
\end{definition}

\begin{definition}
    An arithmetic function is called \textbf{multiplicative} if

    \[f(mn) = f(m)f(n) \ \text{for all coprime} \ m, n.\]
It is called completely multiplicative if

\[f(mn) = f(m)f(n) \ \text{for all} \ m, n.\]
\end{definition}

\begin{definition}
    An arithmetic function is called \textbf{additive} if

    \[f(mn) = f(m)+f(n) \ \text{for all coprime} \ m, n.\]
It is called completely additive if

\[f(mn) = f(m)+f(n) \ \text{for all} \ m, n.\]
\end{definition}

\section{Number of Divisors}

\begin{definition}
Let $n$ be a positive integer. Then the number of divisors of $n$ as a function
is denoted by $d(n)$.    
\end{definition}

\begin{theorem}[Formula for $d(n)$] Let $n \in \mathbb{N}$ such that its prime factorization is

    \[n = p_1^{\alpha_1}p_2^{\alpha_2} \dots p_k^{\alpha_k}\]

Then

    \[d(n) = (1 + \alpha_1)(1 + \alpha_2) \dots (1 + \alpha_k).\]
    
\end{theorem}

\begin{theorem}
    The function $d(n)$ is odd if and only if $n$ is a square.
\end{theorem}

\begin{theorem}
    \[d(n) \leq 2 \sqrt n\]
\end{theorem}

\section{Sum of Divisors}

\begin{definition}
Let $n$ be a positive integer. Then the sum of divisors of $n$ as a function
is denoted by $\sigma(n)$.    
\end{definition}

\begin{theorem}[Formula for $\sigma(n)$] Let $n \in \mathbb{N}$ such that its prime factorization is

    \[n = p_1^{\alpha_1} p_2^{\alpha_2} \dots p_k^{\alpha_k}\]

Then

    \[\sigma(n) = (1 + p_1+\dots+ p_1^{\alpha_1}) \dots (1 + p_k+\dots+p_k ^{\alpha_k}).\]
    
\end{theorem}


\section{Euler's Totient Function}

\begin{definition}
    Let $n$ be a positive integer. The function $\varphi(n)$ is called Euler's totient function, and it denotes the number of positive integers less than $n$ that are coprime to it.
\end{definition}

\begin{theorem}[Gauss] For any positive integer $n$,
    \[\sum_{d\mid n} \varphi(d) = n\]
\end{theorem}

\begin{lemma}
    For all composite $n$, we have
    \[\varphi(n) \leq n-\sqrt{n}\].
    Also, for all $n \neq 2, 6$, 
    \[\varphi(n) \geq \sqrt{n}\]
\end{lemma}

\section{Multiplicative Functions}

\begin{theorem}
    Let $f$ be a multiplicative function. Then 

    \[F(n) = \sum_{d|n} f(d)\]

    is also multiplicative.
\end{theorem}

\begin{lemma}
    $\varphi$ is multiplicative.
\end{lemma}

\begin{definition}
    The \textbf{Möbius function} $\mu(n)$ is defined as follows:   
    \[
        \mu(n) = 
        \begin{cases}
            1 & \text{if } n=1\\
            0 & \text{if } p^2|n \text{ for some prime } p\\
            (-1)^k & \text{if } n=p_1p_2\dots p_k \text{ for distinct primes } p_i
        \end{cases}
    \]
\end{definition}    

\begin{theorem}
   $\mu(n)$ is multiplicative. Also, for any positive integer $n$,
   \[\sum_{d|n} \mu(d) = \begin{cases} 1 & \text{if } n=1 \\ 0 & \text{if } n>1 \end{cases}\]
\end{theorem}

\begin{definition}
    The dirichlet delta function $\delta(n)$ is defined as follows:
    \[\delta(n) = \begin{cases} 1 & \text{if } n=1 \\ 0 & \text{if } n>1 \end{cases}\] 
\end{definition}

    \subsection{Dirichlet Convolution}
        \begin{definition}
            Let $f$ and $g$ be two arithmetic functions (not necessarily multiplicative). Then the Dirichlet convolution of $f$ and $g$ is defined as

            \[(f * g)(n) = \sum_{d|n} f(d)g\left(\frac{n}{d}\right)\]
        \end{definition}

        \begin{theorem}
            The following properties hold for Dirichlet convolution:
            \begin{enumerate}
                \item $f * g = g * f$

                \item $f * (g + h) = f * g + f * h$
                
                \item $(f * g) * h = f * (g * h)$
                
                \item $f * \delta = f$
            \end{enumerate}
        \end{theorem}


        \begin{lemma}
            \[\sigma(n)=\sum_{m\mid n}\varphi(m)\,d\left(\frac{n}{m}\right)\]
        \end{lemma}

    \subsection{Möbius Inversion}
        \begin{theorem}[Möbius Inversion]
            Let $f$ and $g$ be two arithmetic functions such that
            \[g(n) = \sum_{d|n} f(d)\]
            Then
            \[f(n) = \sum_{d|n} g(d)\mu\left(\frac{n}{d}\right)\]

            Equivalently, if $g = f * 1 \ \iff \ f = g * \mu$.
        \end{theorem}
    
\section{Floor and Ceiling Functions}
